% =============================================================================
% Chapter 6: Conclusion -- Part I
% =============================================================================
\mychapter{6. Conclusion}

This mid-term report documents the progress of a master's thesis project on automated XCT defect
detection in additively manufactured metal components using deep learning. The project is on
schedule and has completed the literature review, research design, dataset acquisition and
characterisation, full preprocessing pipeline implementation and validation, patch extraction, data
augmentation, parameter sensitivity analysis, and initial detection model training on the NIST CoCr
\texttt{set1sample2raw} dataset.

\mysection{6.1 Preprocessing Pipeline}
The four-stage preprocessing pipeline, percentile normalisation, beam hardening correction, ring
artefact suppression, and non-local means denoising, was successfully validated on the full NIST CoCr
\texttt{set1sample2raw} dataset (900 slices, $1010\times980$ pixels, raw intensity range
$[0, 59631]$) in 578.5 seconds on CPU. The pipeline achieves a noise standard deviation reduction
from 20952 to 0.385, a factor of $\times 54{,}500$, and an SNR improvement from 1.756 to 1.773.
Visual confirmation (Figures 4.1 and 4.2) demonstrates that defect boundary sharpness is preserved
throughout all preprocessing stages, which is essential for reliable detection of sub-100~$\mu$m
defects.

Parameter sensitivity analysis (Figures 4.6, 4.7, 4.8, 4.9) confirmed the optimal configuration: NLM
filter strength $h = 1.0$ (sweep value; deployed pipeline uses adaptive 0.6), BHC polynomial degree
$d = 3$, and ring filter radius $r = 15$ pixels. This represents an original empirical contribution,
preprocessing parameter selection is frequently omitted or unjustified in published XCT defect
detection studies.

\mysection{6.2 Patch Extraction and Augmentation}
Patch extraction with $256\times256$ pixel patches and 50\% overlap stride was validated on NIST CoCr
data, producing 15,840 patches at $64\times64$ and 3,860 patches at $128\times128$ from 20 sampled
slices (Figures 4.10 and 4.11). The stratified 1:3 foreground-to-background sampling strategy was
implemented and verified to address the severe class imbalance of 1.2--2.8\% defect
voxels~\citep{Sudre2017}.

All six augmentation transforms were validated on NIST CoCr patches (Figure~4.12), confirming correct
spatial alignment between image and label patches across all transform types~\citep{Bellens2021}.

\mysection{6.3 Preliminary Detection Results}
Initial 2D U-Net training on NIST CoCr automatic labels demonstrates clear improvement over both the
classical baseline and the BCE loss formulation:

\begin{table}[htbp]
\centering
\caption{Preliminary detection results after 20 training epochs.}
\label{tab:final_prelim}
\begin{tabular}{lll}
\toprule
\textbf{Method} & \textbf{Validation Dice} & \textbf{vs.\ Target} \\
\midrule
Otsu baseline & 0.38 & --- \\
2D U-Net + BCE & 0.41 & $-45\%$ below target \\
2D U-Net + Dice-Focal & 0.68 & $-9\%$ below target \\
Acceptance criterion & 0.75 & Target \\
\bottomrule
\end{tabular}
\end{table}

The combined Dice-Focal loss already surpasses the Otsu label generator itself (0.68 vs.\ 0.38),
confirming that the U-Net is learning to detect defects beyond the capability of its training
labels~\citep{Yosifov2020}. Full training to convergence and evaluation against the held-out test set
is the immediate next priority.

\mysection{6.4 Key Conclusions}
\begin{itemize}
\item \textbf{Preprocessing is essential and quantifiable}: The four-stage pipeline reduces noise by
  a factor of $\times 54{,}500$ and removes systematic artefacts that would otherwise mask or mimic
  defects. Without preprocessing, reliable automatic label generation and deep learning detection are
  not feasible on this dataset.
\item \textbf{Automatic label generation is viable}: The two-stage Otsu strategy, applying
  thresholding within the metal region only, avoids the background misclassification failure mode of
  single-pass global thresholding. The U-Net already generalises beyond its noisy labels, validating
  this approach~\citep{Yosifov2020,Otsu1979}.
\item \textbf{Loss function choice is critical}: Combined Dice-Focal loss achieves 66\% higher
  validation Dice than BCE (0.68 vs.\ 0.41) after 20 epochs, confirming that class-imbalance-aware
  loss functions are essential for XCT defect detection~\citep{Sudre2017,Lin2017}.
\item \textbf{2D detection is practical and competitive}: The 2D U-Net processes $256\times256$
  patches efficiently, consistent with published findings that 2D approaches achieve within 2--4\% of
  full 3D performance at substantially lower computational cost~\citep{Wong2021}.
\item \textbf{Standardised benchmarking enables reproducibility}: The use of the publicly available
  NIST CoCr AM XCT dataset~\citep{NIST_CoCr} ensures that all results can be independently reproduced
  and compared against future work.
\item \textbf{An end-to-end documented pipeline is the primary deliverable}: The complete pipeline,
  from raw TIFF stack to trained detector with MLflow tracking and DVC dataset versioning, is
  version-controlled on Git and designed for adoption by the AM inspection research community.
\end{itemize}

\mysection{6.5 Remaining Work}
The following work remained after the mid-term submission; Part II of this document reports what has
since been completed:
\begin{itemize}
\item Complete 2D U-Net training to convergence and full evaluation on the held-out NIST CoCr test
  set against all four acceptance criteria (Dice $\geq 0.75$, IoU $\geq 0.60$, Recall $\geq 0.80$ for
  defects $> 50\,\mu$m, inference time $\leq 120$ seconds)
\item Systematic comparison of all four loss function formulations (BCE, Dice, Focal, combined
  Dice-Focal)
\item Independent assessment of each augmentation transform's contribution to detection performance
\item Manual annotation of 5--10 representative NIST CoCr slices for independent ground truth
  validation
\item Application to additional datasets for cross-material generalisation evaluation, completed for
  PODFAM data in Part II
\item Full results, discussion, and conclusion chapter, completed for the PODFAM phase in Part II
\item Final pipeline documentation and open-source release on Git
\end{itemize}
